\chapter{Conclusiones}\label{cha:conclusiones}

	En este capítulo discutimos los principales resultados obtenidos a lo largo de este trabajo y evaluamos el grado de concordancia alcanzado entre el modelo desarrollado y las observaciones disponibles.
	Asimismo, presentamos brevemente algunas perspectivas de trabajo futuro.

	El objetivo central de esta tesis fue investigar los procesos físicos responsables de la emisión no térmica en los \jets asociados a radiogalaxias.
	Con este fin, desarrollamos un modelo teórico que describe el transporte de electrones relativistas a lo largo del \textit{jet}, incorporando tanto los mecanismos de pérdida de energía radiativa y adiabática como el transporte convectivo del plasma.
	A partir de la distribución de partículas obtenida, calculamos la luminosidad sincrotrón producida por los electrones y resolvimos la ecuación de transporte radiativo para un \jet inhomogéneo, lo que permitió obtener predicciones observables directamente comparables con los datos.

	Los resultados presentados en el capítulo anterior indican que el modelo desarrollado es capaz de reproducir de manera consistente las principales características observadas en la emisión de \textit{jets} relativistas asociados a radiogalaxias.
	En particular, para el caso del modelado de Cen~A, se observa que la distribución espectral de energía obtenida presenta la pendiente esperada para una población de electrones relativistas inyectados con un índice espectral $p\simeq2.3$ \eg{en el régimen en que el enfriamiento sincrotrón domina las pérdidas de energía}.

	Asimismo, el modelo reproduce la tendencia observacional del desplazamiento del pico de emisión hacia frecuencias progresivamente más bajas a medida que aumenta la distancia a la base del \textit{jet}.
	Este comportamiento es consistente con la expansión del flujo.
	Como resultado, el plasma se vuelve ópticamente delgado a frecuencias cada vez más bajas conforme se propaga a lo largo del \textit{jet}, en acuerdo con el escenario físico esperado.

	También señalamos que el perfil de temperatura de brillo en función de la altura muestra una buena concordancia global con las observaciones del \textit{jet} de Centaurus A.
	No obstante, es importante destacar que el perfil observado posee una morfología distintiva del modelado, así como diferencias importantes entre ambos "brazos" del \jet.
	Esta discrepancia podría estar asociada a la presencia de perturbaciones o inestabilidades a lo largo del \textit{jet}.
	La inclusión de estos efectos en un modelo más detallado podría contribuir a reproducir con mayor fidelidad la estructura observada del \jet y explicar las diferencias entre ambos lados del flujo.

	Finalmente, el mapa sintético obtenido revela la aparición del fenómeno de bordes brillantes, similares a los observados en Cen A \citep{2021NatAs...5.1017J}, para \jets relativistas.
	En conjunto, estos resultados demuestran que la inclusión de estructura transversal en los modelos de \textit{jets} relativistas constituye un avance significativo respecto de los enfoques unidimensionales tradicionales y ofrece una herramienta útil para la interpretación de observaciones actuales y futuras de radiogalaxias resueltas con técnicas interferométricas.

	En virtud de esto, mencionamos a continuaci\'on algunas posibilidades sobre el futuro de este trabajo:

\begin{itemize}
\item Incorporar la presencia de perturbaciones e inestabilidades dinámicas a lo largo del \textit{jet}, tanto internas al flujo como inducidas por la interacción con el medio.
\item Aplicar el modelo a otras fuentes cercanas con observaciones interferométricas de muy alta resolución angular, tales como M87, con el objetivo de evaluar su capacidad predictiva en distintos regímenes físicos y geométricos.
\item Incluir explícitamente el modelado de la emisión del contra-\textit{jet}, lo que permitiría imponer restricciones más estrictas sobre la geometría, así como los parámetros vinculados con la emisi\'on del flujo, mejorando así el ajuste global a las observaciones.
\item Incorporar la emisión por procesos de Compton inverso (IC) y \textit{Synchrotron Self-Compton} (SSC), con el objetivo de extender el marco desarrollado hacia una descripción multi-frecuencia más completa y comparar con observaciones en bandas de mayor energía.
\end{itemize}


